\section{Ethics Statement}

% 中文注释：伦理说明保持简短克制，只统一正文字体和少量符号间距，不增加新的内容主张。
BiTempQA is intended as an evaluation benchmark for memory systems rather than as a deployed decision-making tool. The dataset is generated from synthetic scenarios and does not intentionally include private personal data. This design reduces direct privacy risk, but it does not eliminate the possibility that generated examples may still reflect cultural, linguistic, or model-induced biases.

% 中文注释：这里保持与 limitations 一致，只去掉多余的样式感。
Our experiments also rely on LLM-generated QA pairs and an LLM-based judge. As discussed in the limitations section, this introduces risks of annotation noise, judge leniency, and benchmark artifacts. We therefore view the benchmark as a diagnostic instrument for controlled comparison rather than as a definitive measure of real-world temporal reliability.

% 中文注释：最后一段继续保持 EMNLP 风格的克制表述。
We recommend that future releases include human validation, broader linguistic coverage, and additional checks for harmful stereotypes or misleading temporal claims. Any use of the benchmark to evaluate production systems in sensitive domains should be accompanied by task-specific safety review and external verification.
